Supervisors of security operations centres must periodically judge whether monitored
organisations detect, investigate, escalate and close security work properly, yet
common instruments are self-assessed maturity models and manual record review. We
describe SAT-SA, a system that analyses the operational records an organisation
submits each period (alerts, cases, investigation steps, escalations, dispositions,
assets) with \V{CODE-workers} deterministic analytical workers, produces findings that
cite those records, ranks entities by risk, requires a human supervisory decision,
and binds the reviewed state with a post-quantum signed, hash-chained receipt. All
results come from frozen, hash-verified experiment bundles. In \V{C01-action-total}
controlled scenarios the system emitted every declared finding family (recall
\V{C01-recall}, precision \V{C01-precision}), but simple fixed-threshold and
median-deviation rules tied its fast-closure detector. Across \V{PR01-replicates}
generated populations its ranking exceeded random and alert-volume ordering, while a
closure-speed heuristic was comparable. Validation matched its declared contract in
all \V{R01b-expect} perturbed conditions, yet stale and contradictory records were
accepted silently. The integrity layer detected all \V{T01-detected} tested
mutations, and all \V{O02-runs} injected workflow interruptions recovered. On a real
public service-desk incident log, which is not security operations data, the
pipeline processed every record, but its risk score was not associated with
service-level misses (Spearman correlation \V{X02b-rho-risk}); a failure analysis
attributes this mainly to construct mismatch. The evidence supports mechanism-level
claims under controlled conditions and feasibility on external data, not
effectiveness in real security operations: no human study or expert labels exist, and
deployment beyond local processes is unverified.
